\documentclass{report}
\usepackage{amsmath,amssymb}
\setlength{\parindent}{0mm}
\setlength{\parskip}{1em}
\begin{document}
\begin{center}
\rule{6in}{1pt} \
{\large Guangye Chen \\
{\bf Nonlinearly implicit, multidimensional, electromagnetic particle-in-cell algorithms for kinetic simulation of plasmas}}

P O Box 1663 Ms K717 \\ Los Alamos National Laboratory \\ Los Alamos \\ NM \\ 87545
\\
{\tt gchen@lanl.gov}\\
Chacon Luis\end{center}

Particle-in-cell (PIC) simulation techniques have been widely successful
in first-principles simulations of plasma dynamics. However, the fundamental
algorithmic underpinnings of standard PIC algorithms have not changed
in decades. The classical PIC method employs an explicit approach
(e.g. leap-frog) to advance the Vlasov-Maxwell/Poisson system using
particles coupled to a grid. Explicit PIC is subject to both temporal
stability constraints (either light-wave CFL condition or resolving
plasma-wave frequency) and spatial stability (so-called finite-grid
instability) constraints, which makes it unsuitable for system-scale
kinetic simulations, even with modern super-computers.

Implicit algorithms can potentially eliminate both spatial and temporal
stability constraints, thus becoming orders of magnitude more efficient
than explicit ones. This has motivated much exploration of these algorithms
in the literature since the 1970's. However, the lack of efficient
nonlinear solvers for a very large system of particle-field equations
required approximations that resulted in intolerable accumulation
of numerical errors in long-term simulations.

In this presentation, we discuss a multi-dimensional, nonlinearly
implicit electromagnetic PIC algorithm \citep{key-1}. The approach
delivers both accuracy and efficiency for multi-scale plasma kinetic
simulations, and extends previous proof-of-principle studies on 1D
electrostatic \citep{key-2,key-3,key-4,key-5} and electromagnetic
systems \citep{key-6}. To avoid noise issues associated with numerical
Cherenkov radiation \citep{key-7}, we focus our implementation on
the Darwin approximation to Maxwell's equations, which eliminates
the light wave analytically. The formulation conserves exactly total
energy, local charge, ignorable canonical momentum, and preserves
the Coulomb gauge. Linear momentum is not exactly conserved, but errors
are controlled by an adaptive particle sub-stepping orbit integrator.
Key to the performance of the algorithm is a moment-based preconditioner,
featuring the correct asymptotic limits. The formulation has been
extended to allow mapped meshes, which opens the possibility of accurate
body-fitted and/or spatially adaptive PIC simulations. The superior
accuracy and efficiency properties of the scheme will be demonstrated
with challenging numerical examples.
\begin{thebibliography}{99}
\bibitem{key-1}G. Chen, L. Chac�n, \textit{Comput. Phys. Commun.}
\textbf{197}, 73-87 (2015).

\bibitem{key-2}G. Chen, L. Chac�n, and D.C. Barnes, \textit{J. Comput.
Phys}., \textbf{230} (18), 7018-7036 (2011).

\bibitem{key-3}G. Chen, L. Chac�n, D. C. Barnes, \textit{J. Comput.
Phys.}, \textbf{231}(16), 5374-5388 (2012).

\bibitem{key-4}L. Chac�n, G. Chen, D. C. Barnes, \textit{J. Comput.
Phys}., \textbf{233}, 1-9 (2013).

\bibitem{key-5}W. Taitano, D. Knoll, L. Chac�n, G. Chen, \textit{SIAM
J. Sci. Compt.}, \textbf{35} (5), S126-S149 (2013).

\bibitem{key-6}G. Chen and L. Chac�n, \textit{Comput. Phys. Commun.},
\textbf{185} (10), 2391-2402 (2014).

\bibitem{key-7}S. Markidis and G. Lapenta, \textit{J. Comput. Phys.},
\textbf{230} (18), 7037-7052 (2011).\end{thebibliography}


\end{document}
